\documentclass[11pt]{article}
\usepackage{fontspec}
\usepackage{amsmath,amssymb}
\usepackage[margin=1in]{geometry}
\usepackage{url}
\begin{document}
\sloppy
\section*{GROK H1602 (2026-10-06 16:02 PT): first-step sweep stopped; 600\,s reruns close 14 of parent cube 0's timed-out subcubes}
\textbf{P versus NP: OPEN.}

\section*{0. What landed since H1545}
\texttt{main} tip is \texttt{25aa6465}, which is the H1545 note itself. No commit landed after it, and the latest commit touching \path{research/main/} is still \texttt{b04263cb} (2026-10-05). Nothing new to read.

\section*{1. Literature}
None. The claim under test is still this loop's own question: is $L(f,3)<9$ for $f=\texttt{a40ff54e580f5d0a}$?

\section*{2. First-step sweep stopped (COMPUTATIONAL, no answer)}
Following the H1545 next step, the uncubed-at-step-2 first-step sweep in \path{/workspace/scratch_h1511_1006/cubes_l8r3.out} was stopped at 16:03\,PT. Its record at stop: cubes 0 to 49 (50 of 150) each timed out at 300\,s; cubes 50 to 53 were killed unanswered; 0 SAT, 0 UNSAT. It produced no information in 51 minutes, so its five cores went to the second-step split.

\section*{3. Parent cube 0, second-step split (COMPUTATIONAL, partial)}
Parent cube 0 is $p=(0,0,1,2)$ with 315 subcubes from the H1531 split (PROVED/EXACT as an enumeration of the tight constraints; positive control in H1545). Read at 16:08\,PT and recounted by script:
\begin{itemize}
\item \path{probe_c0.out} (60\,s cap, still running): subcubes 0 to 181 answered, 148 UNSAT, 34 timed out, 0 SAT.
\item \path{/workspace/scratch_h1600_1006/rerun600.out} (same \path{split2_r3.py}, sha256 \texttt{285c61c4}, 600\,s cap, 5 workers, started 16:03\,PT) reruns the 29 subcubes that had timed out by 16:03. So far 14 returned, all UNSAT, in 66 to 142\,s: 28, 29, 50, 51, 52, 54, 72, 74, 87, 89, 95, 96, 97, 101.
\end{itemize}
Combined: 162 of 315 subcubes UNSAT, 0 SAT; 20 are timed out and not yet resolved (99, 100, 102, 104, 127, 129, 133, 134, 139, 144, 149, 155, 156, 157, 159, 164, 174, 176, 177, 179), and 133 (182 to 314) are not yet run. Parent cube 0 is not closed. Every 60\,s timeout rerun so far has ended UNSAT well under 600\,s, which suggests the 60\,s cap is just short; that is an observation about solver time, not evidence about the remaining subcubes.

\section*{4. Standing status}
No new finite witness improves a charged resource in this pass, and no claim is falsified. UNSAT subcubes cannot improve $T$, $S$, or $R$; they would only support a lower bound once every parent cube is closed, and that bound would still rest on the tightness lemma (CANDIDATE). Current bounds for this $f$, carried from earlier notes and not re-derived here: $L(f,2)\le 10$ (COMPUTATIONAL), $L(f,2)=10$ CONDITIONAL on the tightness lemma; $L(f,3)\le 9$ and $L(f,4)\le 8$ (COMPUTATIONAL, H1444). Whether $L(f,3)=8$ stays OPEN. The time-state-representation surface stays CANDIDATE. The SAT-to-charged-model compiler and the product-span obstruction are both OPEN. The D(4,8) semantic-cut bound of 48 stays FALSIFIED/RETIRED.

\section*{Next smallest step}
Let the 600\,s reruns finish, then rerun the subcubes that timed out after 16:03 (164 onward) at 600\,s. When parent cube 0 has all 315 answered, start parent cube 1 with the 600\,s cap directly. Re-simulate any SAT line with independent code before citing it.

\bigskip
\noindent\textbf{Status sentence.} P versus NP remains OPEN: the SAT-to-charged-model compiler and the product-span obstruction are open load-bearing arrows. This pass stops the uninformative first-step sweep (50 of 150 cubes timed out at 300\,s, 0 answers) and records that parent cube 0 of the tight $L=8$, $r=3$ split for \texttt{a40ff54e580f5d0a} now has 162 of 315 subcubes UNSAT, 20 unresolved timeouts, 133 not yet run, and 0 SAT; no charged resource improves, no claim is falsified, and this is not a proof of anything about P versus NP.
\end{document}
